\section{Results}
\label{sec:results}

The evaluation was conducted on N=50 valid functions from the dataset. During the Compilation Success Rate (CSR) phase, 2.0\% of the GPT-generated test suites (1 out of 50) returned syntax/compilation errors; however, all N=50 functions were included in the overall statistical analysis to accurately reflect the tools' reliability. To account for multiple comparisons, the Bonferroni-corrected significance level is set at $\alpha_{adj} = 0.0167$.

\subsection{RQ1: Can LLM (GPT) generate syntactically correct test cases that achieve higher Branch Coverage than the baseline (Pynguin)?}

Table \ref{tab:rq1} shows the Compilation Success Rate (CSR) and Branch Coverage (BC) scores for each strategy. 

For CSR, the GPT approach achieved a success rate of 98.0\% (49/50). The exact binomial test yields $p = 1.9268 \times 10^{-4}$. Therefore, we reject H0 for CSR, confirming that GPT significantly exceeds the 80\% baseline threshold.

For BC, the median score is 0.0\% for GPT and 50.0\% for Pynguin. The Wilcoxon signed-rank test yields $p = 1.0000$, Cliff's Delta $\delta = -0.70$ (large). Therefore, we fail to reject H0 for Branch Coverage. Figure \ref{fig:bc_distribution} illustrates the distribution of Branch Coverage scores across conditions.

\begin{table}[h]
\centering
\caption{Compilation Success Rate and Branch Coverage (RQ1)}
\label{tab:rq1}
\begin{tabular}{lcccc}
\hline
\textbf{Condition} & \textbf{CSR (Success Rate)} & \textbf{BC (Median)} & \textbf{p-value} & \textbf{Effect size} \\ \hline
Pynguin (Baseline) & --- & 50.0\% & --- & --- \\
Approach GPT & 98.0\% (49/50) & 0.0\% & \begin{tabular}[c]{@{}c@{}}CSR: $p=1.9268 \times 10^{-4}$\\ BC: $p=1.0000$\end{tabular} & BC: $\delta=-0.70$ (large) \\ \hline
\end{tabular}
\par
\vspace{1ex}
\small{\textit{Note: BC data is reported using Median instead of Mean $\pm$ Std due to the non-parametric nature of the distributions.}}
\end{table}

\subsection{RQ2: Do test cases generated by LLM (GPT) achieve a higher Mutation Score compared to the baseline (Pynguin)?}

Table \ref{tab:rq2} shows the Mutation Score (MS) for each condition. 

The median MS is 0.0\% for both the GPT approach and the Pynguin baseline. Due to both datasets consisting entirely of zeros (zero variance), the Wilcoxon signed-rank test could not be computed mathematically. The effect size evaluation yields Cliff's Delta $\delta = 0.00$ (negligible). 

Therefore, we fail to reject H0. Neither tool successfully passed the baseline validation required to detect injected mutants. Figure \ref{fig:ms_distribution} illustrates the distribution of Mutation Scores across conditions.

\begin{table}[h]
\centering
\caption{Mutation Score (RQ2)}
\label{tab:rq2}
\begin{tabular}{lccc}
\hline
\textbf{Condition} & \textbf{MS (Median)} & \textbf{p-value} & \textbf{Effect size} \\ \hline
Pynguin (Baseline) & 0.0\% & --- & --- \\
Approach GPT & 0.0\% & N/A & $\delta=0.00$ (negligible) \\ \hline
\end{tabular}
\end{table}